\section{Continuous maps and sheaves of algebraic structures}
\label{section-presheaves-structures-functorial}

\noindent
Let $(\mathcal{C}, F)$ be a type of algebraic structure.
For a topological space $X$ let us introduce the notation:
\begin{enumerate}
\item $\textit{PSh}(X, \mathcal{C})$ will be the category
of presheaves with values in $\mathcal{C}$.
\item $\Sh(X, \mathcal{C})$ will be the category
of sheaves with values in $\mathcal{C}$.
\end{enumerate}
Let $f : X \to Y$ be a continuous map of topological spaces.
The same arguments as in the previous section show
there are functors
\begin{eqnarray*}
f_* : \textit{PSh}(X, \mathcal{C}) &
\longrightarrow &
\textit{PSh}(Y, \mathcal{C}) \\
f_* : \Sh(X, \mathcal{C}) &
\longrightarrow &
\Sh(Y, \mathcal{C}) \\
f_p : \textit{PSh}(Y, \mathcal{C}) &
\longrightarrow &
\textit{PSh}(X, \mathcal{C}) \\
f^{-1} : \Sh(Y, \mathcal{C}) &
\longrightarrow &
\Sh(X, \mathcal{C})
\end{eqnarray*}
constructed in the same manner and with the same
properties as the functors constructed for abelian
(pre)sheaves. In particular there are commutative diagrams
$$
\xymatrix{
\textit{PSh}(X, \mathcal{C}) \ar[r]^{f_*} \ar[d]^F &
\textit{PSh}(Y, \mathcal{C}) \ar[d]^F &
\Sh(X, \mathcal{C}) \ar[r]^{f_*} \ar[d]^F &
\Sh(Y, \mathcal{C}) \ar[d]^F
\\
\textit{PSh}(X) \ar[r]^{f_*} &
\textit{PSh}(Y) &
\Sh(X) \ar[r]^{f_*} &
\Sh(Y)
\\
\textit{PSh}(Y, \mathcal{C}) \ar[r]^{f_p} \ar[d]^F &
\textit{PSh}(X, \mathcal{C}) \ar[d]^F &
\Sh(Y, \mathcal{C}) \ar[r]^{f^{-1}} \ar[d]^F &
\Sh(X, \mathcal{C}) \ar[d]^F
\\
\textit{PSh}(Y) \ar[r]^{f_p} &
\textit{PSh}(X) &
\Sh(Y) \ar[r]^{f^{-1}} &
\Sh(X)
}
$$

\medskip\noindent
The main formulas to keep in mind are the following
\begin{eqnarray*}
f_*\mathcal{F}(V) & = & \mathcal{F}(f^{-1}(V)) \\
f_p\mathcal{G}(U) & = & \colim_{f(U) \subset V} \mathcal{G}(V) \\
f^{-1}\mathcal{G} & = & (f_p\mathcal{G})^\# \\
(f_p\mathcal{G})_x & = & \mathcal{G}_{f(x)} \\
(f^{-1}\mathcal{G})_x & = & \mathcal{G}_{f(x)}
\end{eqnarray*}
Each of these formulas has the property that they hold in the
category $\mathcal{C}$ and that upon taking underlying
sets we get the corresponding formula for presheaves of sets.
In addition we have the adjointness properties
\begin{eqnarray*}
\Mor_{\textit{PSh}(X, \mathcal{C})}(f_p\mathcal{G}, \mathcal{F})
& = &
\Mor_{\textit{PSh}(Y, \mathcal{C})}(\mathcal{G}, f_*\mathcal{F}) \\
\Mor_{\Sh(X, \mathcal{C})}(f^{-1}\mathcal{G}, \mathcal{F})
& = &
\Mor_{\Sh(Y, \mathcal{C})}(\mathcal{G}, f_*\mathcal{F}).
\end{eqnarray*}
To prove these, the main step is to construct the maps
$$
i_\mathcal{G} : \mathcal{G} \longrightarrow f_*f_p\mathcal{G}
$$
and
$$
c_\mathcal{F} : f_p f_* \mathcal{F} \longrightarrow \mathcal{F}
$$
which occur in the proof of Lemma \ref{lemma-pullback-presheaves} as
morphisms of presheaves with values in $\mathcal{C}$. This may be safely
left to the reader since the constructions are exactly the same
as in the case of presheaves of sets.

\medskip\noindent
Given a continuous map $f : X \to Y$ and sheaves of algebraic
structures $\mathcal{F}$ on $X$, $\mathcal{G}$ on $Y$, the notion
of an {\it $f$-map $\mathcal{G} \to \mathcal{F}$
of sheaves of algebraic structures} makes sense. We can just define
it exactly as in Definition \ref{definition-f-map} (replacing maps
of sets with morphisms in $\mathcal{C}$) or we can
simply say that it is the same as a map of sheaves of algebraic
structures $\mathcal{G} \to f_*\mathcal{F}$. We will use this
notion freely in the following. The set of $f$-maps between
$\mathcal{G}$ and $\mathcal{F}$ will be in canonical bijection
with the sets
$\Mor_{\Sh(X, \mathcal{C})}(f^{-1}\mathcal{G}, \mathcal{F})$
and
$\Mor_{\Sh(Y, \mathcal{C})}(\mathcal{G}, f_*\mathcal{F})$.

\medskip\noindent
Composition of $f$-maps is defined in exactly the
same manner as in the case of $f$-maps of sheaves of
sets. In addition, given an $f$-map $\mathcal{G} \to \mathcal{F}$
as above, the induced maps on stalks
$$
\varphi_x : \mathcal{G}_{f(x)} \longrightarrow \mathcal{F}_x
$$
are homomorphisms of algebraic structures.


\begin{lemma}
\label{lemma-f-map-sets-algebraic-structures}
Let $f : X \to Y$ be a continuous map of topological spaces.
Suppose given sheaves of algebraic structures
$\mathcal{F}$ on $X$, $\mathcal{G}$ on $Y$. Let
$\varphi : \mathcal{G} \to \mathcal{F}$ be an $f$-map
of underlying sheaves of sets. If for every $V \subset Y$ open the
map of sets $\varphi_V : \mathcal{G}(V) \to \mathcal{F}(f^{-1}V)$
is the effect of a morphism in $\mathcal{C}$ on underlying sets,
then $\varphi$ comes from a unique $f$-morphism between
sheaves of algebraic structures.
\end{lemma}

\begin{proof}
Omitted.
\end{proof}
